
%%%%%%%%%% YOU KNOW THE DRILL... REPLACE ALL THE \REPLACEMEs WITH YOUR ANSWERS  %%%%%%%%%

\paragraph{Problem 3 answer}~

\begin{enumerate}[label=(\alph*)]

%%%%%%%%%%%%%%%%%%%%%%%%%%%%%%%%%%%%%%%%%%%%%%%%
%%%%%%%%%%%%%%%%%%%%%%%%%%%%%%%%%%%%%%%%%%%%%%%%

\item                           % (a)

  \(C(1) = \BLUE{
    \REPLACEME 
  }\) 

  \(C(2) = \BLUE{
    \REPLACEME 
  }\)

  \(C(3) = \BLUE{
    \REPLACEME 
  }\)

  \(C(4) = \BLUE{
    \REPLACEME 
  }\) 

%%%%%%%%%%%%%%%%%%%%%%%%%%%%%%%%%%%%%%%%%%%%%%%%
%%%%%%%%%%%%%%%%%%%%%%%%%%%%%%%%%%%%%%%%%%%%%%%%

\item                           % (b)

  \begin{lemma}
    For all $n\ge 1$, executing $C(n)$ does at most $2+2\log_2 n$ additions.
  \end{lemma}

  \begin{longFormProof}

    \step The proof is by induction on $n\ge 1$.

    \step\label{step: base additions}
    For the base case $n=1$, \BLUE{
      \REPLACEME 
    }

    \begin{block}\label{block: mysterious additions}

      \step For the inductive step, consider any $n\ge 2$, and assume the lemma holds for smaller values.

      \begin{blue}

        \step \REPLACEME  
 
        \step \REPLACEME

        % ETC. (YOU MAY NEED MORE THAN TWO STEPS.)

      \end{blue}
      
      \step So $C(n)$ does at most $2 + 2\log_2 n$ additions.
    \end{block}

    \step By Block~\ref{block: mysterious additions}, for any $n\ge 2$, if the lemma holds for smaller values,
    then it holds for $n$.

    \step By induction (with Step~\ref{step: base additions}), the lemma holds for all $n$.
    
  \end{longFormProof}


%%%%%%%%%%%%%%%%%%%%%%%%%%%%%%%%%%%%%%%%%%%%%%%%
%%%%%%%%%%%%%%%%%%%%%%%%%%%%%%%%%%%%%%%%%%%%%%%%

  \continued 

\item \emph{(Problem 3, continued)}   % part (c)
  
  \begin{lemma}
    For all $n\ge 0$, the value output by $C(n)$ is \(\BLUE{
      \REPLACEME
    }\)
  \end{lemma}
  
  \begin{longFormProof}
    
    \step The proof is by induction on $n$.

    \step\label{step: base}
    For the base case $n=0$, \BLUE{
      \REPLACEME
    }

    \begin{block}\label{block: mysterious}

      \step For the inductive step, consider any $n\ge 1$, and assume the lemma holds for smaller values.

      \step Consider the execution of $C(n)$.

      \begin{blue}
        
      \step \REPLACEME

      \step \REPLACEME

      % ETC. (YOU MAY NEED MORE THAN TWO STEPS)

      \end{blue}

      \begin{block}
        \step \underline{Case 1.} First consider the case that $n$ is odd.

        \begin{blue}
          
          \step \REPLACEME

          \step \REPLACEME

          % ETC. (YOU MAY NEED MORE THAN TWO STEPS)

        \end{blue}

      \end{block}

      \begin{block}
        \step \underline{Case 2.} Otherwise $n$ is even. 

        \begin{blue}

          \step \REPLACEME

          \step \REPLACEME

          % ETC. (YOU MAY NEED MORE THAN TWO STEPS) 

        \end{blue}
        
      \end{block}

      \step In either case, $C(n)$ returns \BLUE{\(
        \REPLACEME
        \).}
      That is, the lemma holds for this $n$.
    \end{block}

    \step By Block~\ref{block: mysterious}, for any $n\ge 1$, if the lemma holds for smaller values,
    then it holds for $n$.

    \step By induction (with Step~\ref{step: base}), the lemma holds for all $n$.
    
  \end{longFormProof}

\end{enumerate}


%%%%%%%%%%%%%%%%%%%%%%%%%%%%%%%%%%%%%%%%%%%%%%%%
%%%%%%%%%%%%%%%%%%%%%%%%%%%%%%%%%%%%%%%%%%%%%%%%

\vfill

\paragraph{Problem \arabic{problem} collaboration and citations}~

I collaborated on this problem with the following people:
\begin{blue}
  \begin{itemize}
  \item \REPLACEME                    % write "none" if none
  \end{itemize}
\end{blue}

My answers use ideas from the following sources (other than the lecture notes and textbooks): 
\begin{blue}
  \begin{itemize}
  \item \REPLACEME                    % write "none" if none
  \end{itemize}
\end{blue}


%%% Local Variables:
%%% mode: latex
%%% TeX-master: "homework_1"
%%% End:
